% Section 1: Introduction (~1,500 words)

\subsection{The Problem with Point Estimates}

A redistricting algorithm generates precise numerical outputs: a
Polsby-Popper score of $0.361$, an efficiency gap of $-3.2\%$, a
minority VAP percentage of $51.7\%$ for a particular district. These
numbers appear authoritative. In expert witness testimony, legislative
hearings, or court filings, they carry the weight of algorithmic
objectivity. But what do they actually mean?

A compactness score of $0.361$ is only meaningful in context. If the
algorithm were run with a different random seed, would the score be
$0.347$ or $0.374$? If the Census Bureau had enumerated the population
slightly differently (as it inevitably does, given sampling and coverage
errors), would the district boundaries be different enough to change the
score by $0.010$ or $0.001$? If the underlying geographic units were
defined at block-group rather than tract resolution, would the score
change by $0.020$ or $0.003$?

These are not hypothetical concerns. The Districting Integrity Act
(DIA) \citep{b02paper} requires that algorithmic outputs be certified
as robust to ``documented sources of variation.'' Courts evaluating
algorithmic redistricting claims need to know whether observed
differences between algorithmic and enacted plans exceed the margin of
uncertainty. Expert witnesses who report $+22\%$ compactness improvement
without a confidence interval cannot be cross-examined on whether the
true improvement might be $+5\%$ or $+40\%$.

This paper provides systematic uncertainty quantification for the
METIS recursive bisection redistricting pipeline \citep{b1paper}.
We identify, characterize, and propagate three sources of uncertainty:

\begin{enumerate}
  \item \textbf{Algorithmic seed variance}: METIS uses a random seed
    to initialize its partitioning heuristic; different seeds can
    produce different partitions.

  \item \textbf{Census sampling uncertainty}: The PL 94-171
    redistricting data file \citep{pl94171} is based on a near-complete
    enumeration, but coverage errors and differential undercount create
    uncertainty in district populations.

  \item \textbf{Shapefile resolution measurement error}: The geographic
    boundaries used to construct the adjacency graph are derived from
    TIGER/Line shapefiles \citep{tiger-census}, which approximate true
    boundaries at a finite resolution. Different resolutions produce
    different perimeter and area measurements, and therefore different
    compactness scores.
\end{enumerate}

\subsection{Why UQ Matters for Litigation}

The post-\textit{Rucho} landscape of redistricting litigation places
algorithmic redistricting plans in a distinctive legal role: as
\emph{baseline comparators} \citep{rucho2019}. State courts, following
the reasoning of \textit{LWV v.\ Pennsylvania} \citep{lwv2018} and
\textit{Harper v.\ Hall} \citep{harper2022}, have increasingly accepted
algorithmic plans as evidence that enacted plans could have been drawn
less unfairly. When an expert witness testifies that ``the enacted plan
has an efficiency gap of $+7.5\%$ while the algorithmic baseline is
$-3.2\%$, a difference of $10.7$ percentage points,'' opposing counsel
will ask: \textit{what is the margin of error on that difference?}

Without UQ, the answer is ``I don't know.'' With UQ, the answer is
``The 95\% CI for the difference is $[8.1, 13.3]$ percentage points,
so the direction and approximate magnitude are robust to all documented
sources of variation.'' The latter answer is both scientifically honest
and legally more compelling.

\subsection{Overview of UQ Methods Used}

We employ three statistical methods, matched to the nature of each
uncertainty source:

\textbf{Bootstrap confidence intervals} \citep{efron1993introduction}
for seed variance and partisan metrics. We treat each of the 10,000
seed realizations from B.7 \citep{b7paper} as a bootstrap sample,
enabling direct empirical estimation of the sampling distribution of
any metric without distributional assumptions.

\textbf{Delta method} \citep{ver2012numerical} for census sampling
uncertainty. When the input quantities (population counts) have known
or estimable standard errors, the delta method propagates those errors
to derived quantities (deviation from ideal district population) using
first-order Taylor expansion. This is computationally efficient and
exact for smooth, differentiable metrics like deviation from ideal.

\textbf{Sensitivity analysis} for resolution measurement error.
Since resolution is not a random variable (it is an analytic choice),
we quantify measurement error by comparing metric values across the
three resolutions tested in C.1 \citep{dellalibera2026c1}, treating
the block-level results as the ``truth'' and computing signed error
for tract- and block-group-level measurements.

\subsection{Key Contributions}

This paper makes four contributions:

\begin{enumerate}
  \item A complete CI table for Polsby-Popper compactness, covering
    all three uncertainty sources, for all 50 states at 2020 Census.

  \item Bootstrap CIs for efficiency gap, mean-median difference, and
    partisan bias at 50\% vote share for 15 competitive states.

  \item Delta method CIs for district population deviation from ideal,
    propagating PL 94-171 census undercount estimates from
    \citet{brown2018demographic} and \citet{daponte2021census}.

  \item A synthesized CI for the headline ``$+22\%$ compactness
    improvement'' that accounts for all three sources simultaneously
    using variance decomposition.
\end{enumerate}

\subsection{Paper Structure}

Section~\ref{sec:seed} characterizes METIS seed variance using the
B.7 dataset. Section~\ref{sec:census} derives delta method CIs for
population deviation using published census undercount estimates.
Section~\ref{sec:pp} characterizes PP measurement error from resolution
choice. Section~\ref{sec:partisan} presents bootstrap CIs for partisan
fairness metrics. Section~\ref{sec:vra} addresses uncertainty in
minority VAP estimates. Section~\ref{sec:synthesis} constructs the
joint CI for the $+22\%$ compactness improvement. Section~\ref{sec:conclusion}
concludes with practical guidance for expert witnesses.
